\section{Shortest paths}

\entry{Dijkstra}{$O((n + m) \log n)$ \quad mem $n + m$}{%
\textbf{Non-negative} weights only --- one negative edge and the answer is silently
wrong, use \code{bellman}. \code{d[v] == INF} means unreachable. Pass \code{\&par} to
recover a path: walk \code{par} back from the target and reverse.\\[0.3ex]
Variants that are the same loop with a different relax step: maximise the bottleneck
(\code{max(d[u], w)} with a max-heap), count shortest paths (a second array,
\code{cnt[v] += cnt[u]} on tie), fewest edges among shortest paths (compare
\code{pair<dist, edges>}). For 0/1 weights use a \code{deque} instead of the heap
(push\_front on 0) and drop the log.}

%LABEL dijkstra
\begin{lstlisting}
// g[u] = {(v, w)}
vi dijkstra(const vector<vector<pii>>& g, int s,
            vector<int>* par = nullptr) {
    int n = (int)g.size();
    vi d(n, INF);
    if (par) par->assign(n, -1);
    priority_queue<pii, vector<pii>, greater<pii>> q;
    d[s] = 0; q.push({0, s});
    while (!q.empty()) {
        auto [dv, u] = q.top(); q.pop();
        if (dv > d[u]) continue;            // stale entry
        for (auto [v, w] : g[u])
            if (dv + w < d[v]) {
                d[v] = dv + w;
                if (par) (*par)[v] = (int)u;
                q.push({d[v], v});
            }
    }
    return d;
}
\end{lstlisting}

\entry{Bellman-Ford}{$O(nm)$ \quad mem $n + m$}{%
Handles negative weights and reports a negative cycle reachable from \code{s}. Stops
early when a round changes nothing, so it is fast on easy inputs.\\[0.3ex]
\textbf{Difference constraints.} A system of $x_v - x_u \le w$ is exactly a shortest
path problem: add edge \code{u -> v} of weight \code{w}, add a super-source with a
0-edge to every vertex, run this from it. The resulting \code{d} is a feasible
assignment, and the system is infeasible iff there is a negative cycle. $x_v - x_u \ge
w$ becomes $x_u - x_v \le -w$; $x_v - x_u = w$ is both.\\[0.3ex]
To mark every vertex whose distance is truly $-\infty$, run $n$ more rounds and set
\code{d[v] = -INF} whenever it still relaxes.}

\begin{lstlisting}
// e = {u, v, w}; neg = a negative cycle is reachable from s
vi bellman(int n, const vector<array<ll,3>>& e, int s,
           bool& neg) {
    vi d(n, INF); d[s] = 0; neg = false;
    rep(i, 0, n) {
        bool ch = false;
        for (auto [u, v, w] : e)
            if (d[u] < INF && d[u] + w < d[v]) {
                d[v] = d[u] + w; ch = true;
                if (i == n - 1) neg = true;
            }
        if (!ch) break;
    }
    return d;
}
\end{lstlisting}

\entry{Floyd-Warshall}{$O(n^3)$ \quad mem $n^2$}{%
All pairs, negative weights allowed. Fill \code{d[i][i] = 0}, \code{d[i][j]} = the
\emph{cheapest} edge \code{i -> j} (careful with multi-edges: take a \code{min}), and
\code{INF} elsewhere. Rewrites \code{d} in place.\\[0.3ex]
A negative cycle exists iff \code{d[i][i] < 0} for some \code{i} afterwards; then
\code{d[a][b]} is $-\infty$ whenever \code{d[a][i] < INF} and \code{d[i][b] < INF} for
such an \code{i}. The \code{k} loop must be outermost. The \code{< INF} guards are what
stop \code{INF + INF} from wrapping.\\[0.3ex]
Same triple loop, other semirings: transitive closure (\code{|=} with \code{\&\&}),
maximum bottleneck (\code{max}/\code{min}), minimax path.}

\begin{lstlisting}
void floyd(matrix& d) {
    int n = (int)d.size();
    rep(k, 0, n) rep(i, 0, n) if (d[i][k] < INF)
        rep(j, 0, n) if (d[k][j] < INF)
            d[i][j] = min(d[i][j], d[i][k] + d[k][j]);
}
\end{lstlisting}

\entry{K shortest walks}{$O((m + k n)\log)$ \quad mem $n + m$}{%
The \code{k} smallest \code{s}$\to$\code{t} lengths in ascending order, \textbf{walks}
--- vertices may repeat, so on a graph with a cycle you get the same path plus loops.
That is what most ``k-th shortest route'' problems want. Non-negative weights.
Fewer than \code{k} results means fewer exist.\\[0.3ex]
For genuinely \emph{simple} k shortest paths you need Yen's algorithm --- much longer,
and rarely what is asked.}

\begin{lstlisting}
vi kShortest(const vector<vector<pii>>& g, int s, int t,
             int k) {
    int n = (int)g.size();
    vector<int> cnt(n, 0); vi res;
    priority_queue<pii, vector<pii>, greater<pii>> q;
    q.push({0, s});
    while (!q.empty() && (int)res.size() < k) {
        auto [d, u] = q.top(); q.pop();
        if (cnt[u] >= k) continue;
        cnt[u]++;
        if (u == t) res.push_back(d);
        for (auto [v, w] : g[u])
            if (cnt[v] < k) q.push({d + w, v});
    }
    return res;
}
\end{lstlisting}
